\section{Global limits and local conservation}\label{sec:compactness}

The uniform bounds of Proposition~\ref{prop:global-control} allow us to
remove both the collision cutoff and the truncation of the initial jets.
The relevant compactness is stronger than a bound on mass and entropy:
the same-position collision products are uniformly integrable, with
polynomial velocity weights. This will preserve the local energy flux
and rule out defects in all five collision-invariant balances.
The argument extends the fixed-horizon compactness proof of
\cite[Section~6]{OpenAIPeriodic} to the global approximations constructed
here; we give the proof, including the averaging step.

\begin{proposition}\label{prop:global-limits}
Fix the construction parameters supplied by
Proposition~\ref{prop:global-control}, and let $b_N\to\infty$.
Either the dormant family, or a family of cap solutions at one fixed horizon
$T\in(0,T_0]$, has a subsequence of total densities $F_N$ converging
weakly in $L^1([0,\tau]\times\Omega)$ for every $\tau<\infty$ to a
nonnegative density $F$. This subsequence can be chosen so that
$F_N(t)\rightharpoonup F(t)$ in $L^1(\Omega)$ at every $t\ge0$,
with uniformly small polynomially weighted velocity tails on each
bounded time interval.

The limit has initial datum $F_0$ and is a global entropy solution
obeying local conservation in the sense of
Definition~\ref{def:admissible}. More precisely,
\begin{equation}\label{comp:strong-properties}
 F\in C([0,\infty);L^1(\Omega)),\qquad
 \int_0^\tau\!\int_\Omega
 \langle v\rangle^k\bigl(F+Q^+(F,F)+Q^-(F,F)\bigr)
 \dd x\dd v\dd t<\infty
 \quad(k\ge0,\ \tau<\infty).
\end{equation}
Total mass, momentum, and energy are conserved at every time.
The common early-time pointwise majorant of the dormant approximations
also bounds any dormant limit.
\end{proposition}

Throughout the proof, $T$ denotes the fixed cap horizon when the cap
family is used, whereas $\tau$ denotes an arbitrary finite observation
time. Constants may depend on $\tau$, but not on $N$ or $b_N$.
We write
\[
 Q_N^\pm=Q^{\pm,b_N}(F_N,F_N),\qquad
 R_{k,N}(t,x)=\int_{\R^3}\langle v\rangle^kF_N(t,x,v)\dd v.
\]
The initial convergence is supplied by Lemma~\ref{con:datum};
the finite-approximation entropy inequalities and local conservation
laws are those of Proposition~\ref{con:finite-evolution}.

\subsection{Small scales and uniform integrability}

Recall that a color of scale $s=s_j$ remains active until
$L_s=s^{-1/2}$. Let $\mathcal A_s(t)$ be the union, in one periodic
cell, of its two signs' enlarged sets $\mathcal O_i(t)$, including
all lattice copies, when $t\le L_s$; set
$\mathcal A_s(t)=\varnothing$ for $t>L_s$.
We use these sets for every scale, even if that scale is absent from a
particular finite approximation. The active enlarged sets are disjoint.
Their axial lengths on $[0,\tau]$ are bounded by $C_\tau$, so
\[
 |\mathcal A_s(t)|\le C_\tau s^2
 \qquad(0\le t\le\tau).
\]
The free-core amplitude estimate, the cold remainder bound, and the
bounds on the bath and the hot or regular part imply
\begin{equation}\label{comp:moment-envelope}
 R_{k,N}(t,x)\le C_{k,\tau}
 \left(1+\sum_{j\ge1}
       \frac{\ind_{\mathcal A_{s_j}(t)}(x)}{t+s_j}\right).
\end{equation}
Here the bounded term includes all nonmain colors and all retired
colors: after retirement they belong to $w\le BH$.
For active main colors the estimate follows from
Lemma~\ref{cold:core-amplitude} and Proposition~\ref{cold:control};
the nonmain contribution is controlled by Lemma~\ref{cold:nonmain}.

The square of this envelope controls the collision products.
Disjointness ensures that no product of two distinct singular terms
appears when it is squared, and one scale contributes at most
\begin{equation}\label{comp:annular-square}
 \int_0^\tau\frac{C_\tau s^2}{(t+s)^2}\dd t
 =C_\tau s^2\left(\frac1s-\frac1{\tau+s}\right)
 \le C_\tau s.
\end{equation}
For $\varepsilon>0$, put
$\mathcal A_{<\varepsilon}(t)=
\bigcup_{s_j<\varepsilon}\mathcal A_{s_j}(t)$.
Since the scales form a geometric sequence,
\begin{equation}\label{comp:small-annuli}
 \sup_N\int_0^\tau\!\int_{\mathcal A_{<\varepsilon}(t)}
          R_{k,N}(t,x)^2\dd x\dd t
 \le C_{k,\tau}\sum_{s_j<\varepsilon}(s_j+s_j^2)
 \longrightarrow0
 \qquad(\varepsilon\downarrow0).
\end{equation}
The first power of the envelope gives, uniformly at individual times,
\begin{equation}\label{comp:fixed-time-small}
 \sup_{N,\,0\le t\le\tau}
 \int_{\mathcal A_{<\varepsilon}(t)}R_{k,N}(t,x)\dd x
 \le C_{k,\tau}\sum_{s_j<\varepsilon}(s_j+s_j^2)
 \longrightarrow0.
\end{equation}
These estimates include retirement times, because a color may be
included in its last active enlargement at that section and is covered
by the bound on $w$ thereafter.

Outside $\mathcal A_{<\varepsilon}(t)$, a main color can have only
one of the finitely many scales $s_j\ge\varepsilon$.
Their thermal widths are positive and fixed. The Gaussian velocity
bounds and the summation over lattice copies therefore give
\begin{equation}\label{comp:fixed-majorant}
 F_N(t,x,v)\le \Psi_{\varepsilon,\tau}(v)
 \quad\text{if }x\notin\mathcal A_{<\varepsilon}(t),
 \qquad
 0\le\Psi_{\varepsilon,\tau}\le C_{\varepsilon,\tau}H.
\end{equation}
The same majorant includes the bounded hot, bath, and retired parts,
and the Gaussian nonmain tails. In particular it has every polynomial
velocity moment; its singularity on the velocity axis is integrable.

Equations~\eqref{comp:small-annuli}--\eqref{comp:fixed-majorant}
prove uniform integrability in the full spacetime and two-velocity
variables of
\begin{equation}\label{comp:weighted-pairs}
 \langle v\rangle^k\langle v_*\rangle^k
 F_N(t,x,v)F_N(t,x,v_*),\qquad k\ge0.
\end{equation}
Indeed their integrals on the removed spatial sets are bounded by
\eqref{comp:small-annuli}; on the complement they are dominated by the
fixed integrable function
$\langle v\rangle^k\Psi_{\varepsilon,\tau}(v)
 \langle v_*\rangle^k\Psi_{\varepsilon,\tau}(v_*)$.
First make the removed contribution small, and then use absolute
continuity of this fixed integral on arbitrary small measurable sets.
The same decomposition gives uniformly small weighted velocity tails.
Using the first power instead of the square proves the corresponding
one-particle assertions, both in spacetime and, by
\eqref{comp:fixed-time-small}, uniformly at all individual times.

The separate collision terms satisfy the same spacetime conclusions.
Collision symmetry and conservation of the sum of squared speeds give
\begin{equation}\label{comp:collision-moments}
 \int_{\R^3}\langle v\rangle^kQ_N^\pm(t,x,v)\dd v
 \le C_k R_{k+1,N}(t,x)^2.
\end{equation}
For the gain, this follows by placing the velocity weight on an
outgoing particle and using
$\langle v'\rangle^k|v-v_*|
 \le C_k\langle v\rangle^{k+1}\langle v_*\rangle^{k+1}$;
the loss is estimated directly. The cutoff factor is at most one.
Thus \eqref{comp:small-annuli} controls their weighted integrals on the
removed sets. On the complement, positivity gives
\[
 Q_N^+\le P(\Psi_{\varepsilon,\tau},\Psi_{\varepsilon,\tau}),
 \qquad
 Q_N^-\le\Psi_{\varepsilon,\tau}\nu_{\Psi_{\varepsilon,\tau}}.
\]
Both majorants have every polynomial moment, either by
\eqref{comp:collision-moments} for the fixed density
$\Psi_{\varepsilon,\tau}$ or by the bounds on $H$.
Consequently $\{Q_N^\pm\}$ is uniformly integrable with every fixed
polynomial velocity weight. Its integrals over time intervals whose
length tends to zero vanish uniformly in $N$: the small-scale
contribution can first be discarded, and the remaining fixed
majorants have integrals proportional to the length of the interval.

\subsection{Velocity averaging with uniformly integrable sources}

Weak compactness alone does not identify a product evaluated at the
same spatial point. We obtain the needed compactness of velocity
averages from the transport equation. The following elementary form of
the averaging principle of Golse, Lions, Perthame, and Sentis
\cite{GolseLionsPerthameSentis1988} is also proved in
\cite[Section~6]{OpenAIPeriodic}.

\begin{lemma}\label{comp:averaging}
Let $f_n,g_n$ be locally uniformly integrable in $L^1$ on an open
region of time, space, and velocity, and suppose that
$(\partial_t+v\cdot\nabla_x)f_n=g_n$ there in distributions.
For every smooth compactly supported velocity test $\psi$, the averages
$\int\psi(v)f_n(t,x,v)\dd v$ are relatively compact in local
$L^1_{t,x}$ wherever the hypotheses hold for all velocities in a
neighborhood of $\supp\psi$.
The spatial domain may be a flat torus or an open subset of $\R^3$.
\end{lemma}

\begin{proof}
Localize to a compact time-space set and a compact velocity
neighborhood of $\supp\psi$, using smooth cutoffs. Denote the localized
functions by $q_n$, extend them by zero, and set
$r_n=(1+\partial_t+v\cdot\nabla_x)q_n$.
The equation shows that $r_n$ is uniformly integrable in $L^1$ and has
support in a fixed bounded set of time and velocity, and also of space
in the Euclidean case. Spatial cutoff derivatives cause no difficulty
because velocity has been localized. Integration from a time preceding
the support gives
\[
 q_n(t,x,v)=\int_0^\infty e^{-s}r_n(t-s,x-sv,v)\dd s=:Rr_n(t,x,v).
\]
The damped inverse $R$ is an $L^1$ contraction, since free transport
preserves Lebesgue measure.

Truncate $r_n$ at height $L$, preserving its sign, to obtain
$r_n^{(L)}$. Uniform integrability implies
\[
 \sup_n\|r_n-r_n^{(L)}\|_1\longrightarrow0
 \quad(L\to\infty),\qquad
 \sup_n\|r_n^{(L)}\|_2<\infty
 \quad\text{for fixed }L.
\]
The velocity averages of $R(r_n-r_n^{(L)})$ are therefore uniformly
small in $L^1$. It remains to prove compactness for a bounded $L^2$
family of right-hand sides.

Take Fourier transform in $(t,x)$, with Fourier series in the torus
case. For $u_n=Rr_n^{(L)}$, Cauchy--Schwarz yields
\[
 \left|\int\psi(v)\widehat u_n(\sigma,k,v)\dd v\right|^2
 \le I_\psi(\sigma,k)
       \int|\widehat r_n^{(L)}(\sigma,k,v)|^2\dd v,
 \qquad
 I_\psi(\sigma,k)=\int
 \frac{|\psi(v)|^2}{1+(\sigma+k\cdot v)^2}\dd v.
\]
If $\supp\psi\subset\{|v|\le R_0\}$, integration in the direction
of $k$ gives $I_\psi\le C_\psi/|k|$ for $k\ne0$.
When $|\sigma|>2R_0|k|$, the denominator instead gives
$I_\psi\le C_\psi/(1+|\sigma|^2)$.
Thus $I_\psi\to0$ uniformly as $|\sigma|+|k|\to\infty$.
Plancherel's theorem shows that the averages have uniformly small
high-frequency $L^2$ tails. Convolution with a fixed smooth approximate
identity consequently approximates them uniformly in $L^2$.
On every compact time-space set these convolutions are bounded and
equicontinuous, by Cauchy--Schwarz applied to translates of the
convolution kernel and its derivatives. They are relatively compact
there. This proves local $L^2$, and hence local $L^1$, compactness at
fixed $L$; the uniform $L^1$ error as $L\to\infty$ completes the proof.
The argument applies unchanged to the dual lattice of a rotated torus.
\end{proof}

\subsection{Identification of the collision terms}

The uniform integrability and velocity tightness give a subsequence
with $F_N\rightharpoonup F$ in
$L^1([0,\tau]\times\Omega)$.
Since $\D F_N=Q_N^+-Q_N^-$, Lemma~\ref{comp:averaging} applies.
After extraction for a countable dense family of compact velocity
tests, all smooth compact velocity averages converge strongly in local
$L^1_{t,x}$ to the corresponding averages of $F$.
The extension from the countable family follows from the uniform
integrated mass bound.

Set $W_N(t,x,v,v_*)=F_N(t,x,v)F_N(t,x,v_*)$.
By \eqref{comp:weighted-pairs}, a further subsequence has a weak $L^1$
limit $W$. For smooth compactly supported $\psi,\chi$,
\[
 \int\psi(v)\chi(v_*)W_N(t,x,v,v_*)\dd v\dd v_*
 =\left(\int\psi F_N\dd v\right)
  \left(\int\chi F_N\dd v_*\right).
\]
The right-hand side converges in measure locally in $(t,x)$.
It is uniformly integrable there: its absolute value is bounded by the
integral of $|\psi(v)\chi(v_*)|W_N$, and the full pair-product bounds
control this integral on small time-space sets. Vitali's theorem
therefore gives convergence in local $L^1$.
It follows that $W$ and $F(t,x,v)F(t,x,v_*)$ have identical integrals
against separated velocity tests, for almost every $(t,x)$.
A countable determining family gives a common full-measure set;
finite sums of separated tests are dense among continuous functions on
compact two-velocity sets. Hence
\begin{equation}\label{comp:pair-limit}
 W_N\rightharpoonup F(t,x,v)F(t,x,v_*)
 \quad\text{in }L^1([0,\tau]\times\Omega\times\R^3).
\end{equation}
In particular the product on the right is integrable.
The uniform integrability removes time strips near $0$ and $\tau$,
so interior averaging suffices. Weighted tightness extends the
identification to each fixed polynomial velocity weight. One may
extract simultaneously for integer moment orders; domination by a
higher integer order then supplies every order $k\ge0$.

Adjoining the finite angular measure preserves weak convergence.
For each $\omega$, the elastic substitution on $(v,v_*)$ is
measure preserving and involutive. Consequently
\eqref{comp:pair-limit} also gives weak convergence of the primed pair
products on the collision space. On bounded input velocities the
collision factor is bounded and the cutoff equals one for all large
$N$. The estimates with higher velocity moments remove the complement.
For the gain this uses the preserved sum of squared speeds when an
output test is pulled back to the inputs. These observations identify
any weak limit of the separately uniformly integrable gains and losses:
\begin{equation}\label{comp:collision-limit}
 Q_N^\pm\rightharpoonup Q^\pm(F,F)
 \quad\text{in }L^1([0,\tau]\times\Omega),
\end{equation}
also after multiplication by any fixed polynomial velocity weight.
The argument permits bounded measurable tests: after the angular
integration and elastic substitution, such a test remains bounded on
each truncated collision domain.

We can now pass to the unrenormalized equation, including the initial
boundary term, since the initial densities converge in $L^1$.
Thus $\D F=Q(F,F)$ in distributions with datum $F_0$, and both
separate collision terms have the weighted integrability in
\eqref{comp:strong-properties}. The remaining work concerns the time
representative, entropy, and the local invariant balances.

\subsection{A single representative at every time}

For a smooth periodic $\varphi(x,v)$ compactly supported in velocity,
the cutoff path equations give
\[
 \big|\langle F_N(t)-F_N(s),\varphi\rangle\big|
 \le\int_s^t\!\int_\Omega
 \left(F_N|v\cdot\nabla_x\varphi|
       +(Q_N^++Q_N^-)|\varphi|\right)\dd x\dd v\dd u.
\]
The first term is bounded by $C_\varphi|t-s|$.
The second has a modulus tending to zero with $|t-s|$, uniformly in
$N$, by the short-time-integral conclusion above.
For a countable family dense in the uniform norm on each compact
velocity set, Arzel\`a--Ascoli therefore supplies uniform convergence
of the pairings on $[0,\tau]$, after one extraction.

At any individual time, \eqref{comp:fixed-time-small} and
\eqref{comp:fixed-majorant} make $\{F_N(t)\}$ relatively weakly
compact in $L^1(\Omega)$. Every cluster point has the smooth pairings
just obtained, and hence is the same density. Thus this one subsequence
converges weakly in $L^1$ at every time, defining a representative
$F(t)$ whose smooth pairings are continuous and which agrees with the
spacetime limit. Weighted fixed-time tightness gives the same conclusion
for moments, so the conserved cutoff mass, momentum, and energy pass
at every time. In particular the initial trace of this representative
is $F_0$.

Let $\Phi(A,B)=(A-B)\log(A/B)$ with its nonnegative
lower-semicontinuous extension. It is jointly convex: it is the sum of
the two convex perspectives $A\log(A/B)$ and $B\log(B/A)$.
On any bounded collision domain, \eqref{comp:pair-limit} and its primed
version give joint weak convergence of the two arguments of $\Phi$.
Restriction to $[0,t]$ preserves this convergence for every fixed
$t\in[0,\tau]$. Weak lower semicontinuity therefore applies to the
integral of $|a|\Phi$; the weight can be absorbed into the arguments
by $\Phi(cA,cB)=c\Phi(A,B)$ for $c\ge0$.
On $|v|,|v_*|\le R$ the cutoff is absent once $b_N\ge2R$.
Exhausting the velocity space yields
\[
 \int_0^t\mathcal D(F(s))\dd s
 \le\liminf_N\int_0^t\mathcal D^{b_N}(F_N(s))\dd s.
\]
The same convex lower-semicontinuity argument applies to
$\mathcal H_M(F_N(t))$ under the weak $L^1$ convergence at that time.
Combining the two nonnegative terms with the finite-approximation
entropy inequality and initial entropy convergence gives
\begin{equation}\label{comp:limit-entropy}
 \mathcal H_M(F(t))+\int_0^t\mathcal D(F(s))\dd s
 \le\mathcal H_M(F_0),\qquad 0\le t\le\tau.
\end{equation}
The extraction is independent of $t$; no exceptional set of time
endpoints is introduced by the lower-semicontinuity argument.

For completeness, the negative part of $F\log F$ is bounded by mass,
energy, and a fixed Gaussian integral. On
$e^{-1-|v|^2}\le F<1$, use
$-F\log F\le(1+|v|^2)F$; below that threshold, monotonicity of
$-r\log r$ on $(0,e^{-1})$ gives the integrable bound
$(1+|v|^2)e^{-1-|v|^2}$.
The positive part is controlled by the relative entropy integrand and
a multiple of $F$. Hence the absolute logarithmic integral is bounded
uniformly on $[0,\tau]$.

\subsection{Strong traces, renormalization, and the invariant balances}

The integrability of the unrenormalized collision operator now yields
strong continuity. In free coordinates the distributional equation,
including its initial trace, reads
\[
 \partial_t\bigl[F(t,x+tv,v)\bigr]=Q(F,F)(t,x+tv,v).
\]
The right-hand side belongs to spacetime $L^1$.
Subtracting its Bochner integral leaves a distribution with zero time
derivative and initial value $F_0$. Therefore
\begin{equation}\label{comp:path-formula}
 F(t,x+tv,v)=F_0(x,v)
       +\int_0^t Q(F,F)(s,x+sv,v)\dd s
\end{equation}
holds as an $L^1_{x,v}$ identity, first for almost every time.
Its right-hand side is absolutely continuous in $L^1$.
Free transport is a strongly continuous group on $L^1(\Omega)$,
as follows by approximation by continuous functions with compact
velocity support. Thus \eqref{comp:path-formula} defines a strongly
continuous representative. Its smooth pairings agree almost
everywhere in time with those of the representative already chosen;
continuity makes them agree at every time.
In particular $F\in C([0,\tau];L^1)$, $F(0)=F_0$ strongly, and
pairings with every bounded measurable test are continuous.
All these observations also hold with any fixed polynomial velocity
weight, since both the datum and the source have that integrability.

Equation~\eqref{comp:path-formula} holds on almost every individual
free path. On such a path $M(v)>0$ is constant. For a renormalization
$\beta$ from Definition~\ref{def:admissible}, the ordinary absolutely
continuous chain rule gives
\[
 \D\bigl[M\beta(F/M)\bigr]=\beta'(F/M)Q(F,F).
\]
The source is integrable because $|\beta'|\le C_\beta$.
Moreover, for nonnegative $f,g$,
\[
 |M\beta(f/M)-M\beta(g/M)|\le C_\beta|f-g|,
 \qquad |M\beta(f/M)|\le C_\beta f.
\]
The strong initial and terminal traces consequently pass through the
renormalization. Integrating the path identity against a smooth test
proves exactly \eqref{eq:renormalized}, including its initial term.
The separate integrability requirement follows from
$0\le Q^\pm/(1+F/M)\le Q^\pm$.

Finally, let
$\psi\in\{1,v_1,v_2,v_3,|v|^2/2\}$ and let
$\phi\in C_c^\infty([0,\infty)\times(\R^3/\Lambda))$.
At the cutoff level the local identity is
\[
 \int_0^\infty\!\int_\Omega
 F_N\psi(v)(\partial_t\phi+v\cdot\nabla_x\phi)
 \dd x\dd v\dd t
 +\int_\Omega F_N(0,x,v)\psi(v)\phi(0,x)\dd x\dd v=0.
\]
It follows from the path equation and collision symmetry, all integrals
being absolutely convergent. Weighted weak convergence of $F_N$ and
weighted convergence of the initial data pass this identity to $F$.
For the energy identity the largest required weight is $|v|^3$,
which is covered by \eqref{comp:strong-properties}.
This proves all the exact local balances \eqref{eq:local-balances}
with their initial traces. Equivalently, their collision terms vanish
in the limit by absolute weighted collision integrability; no
moment defect is present.

Apply the preceding extractions successively on $[0,m]$, $m\in\N$,
and take a diagonal subsequence, retaining the countable moment and
velocity-test families. The representatives on overlapping intervals
agree by their smooth pairings, so they define one global strongly
continuous density. All integrability, entropy, and conservation
statements hold on every finite interval. For a dormant sequence, sum
the cold majorants over all scales and add the bath and dormant-hot
majorants on their common early interval. This is a locally integrable
pointwise upper bound independent of $N$. Testing against nonnegative
compactly supported functions and passing to the weak limit preserves
it. Its velocity integral is also integrable in space at each fixed
time, by the first-power estimate used in
\eqref{comp:fixed-time-small}. The every-time weak convergence therefore
preserves the bound at each such time section as well. This completes
the proof of Proposition~\ref{prop:global-limits}.
